% ============================================================
% Chương 8. Tổng kết
% Nguồn: docs/report_tv1.md mục 7, 8
% ============================================================
\section{Tổng kết}
\label{sec:tong-ket}

\subsection{Kết quả đạt được}

Nhóm đã xây dựng một kho dữ liệu phân tích học tập hoàn chỉnh từ bộ dữ liệu OULAD, đi từ 7
file CSV gốc đến dashboard phục vụ bốn nhóm người dùng. Về thiết kế, dữ liệu được phân tích
qua ba mức ý niệm, luận lý và vật lý; các phụ thuộc hàm được kiểm chứng trên dữ liệu thật để
đưa lược đồ về dạng chuẩn, rồi phi chuẩn hóa có chủ đích thành dimensional model gồm 5
dimension và 4 fact. Cùng bộ dữ liệu còn được tổ chức theo các mô hình 3NF, wide table, dạng
lồng và tài liệu, kèm số đo cho thấy mỗi mô hình phù hợp với kiểu truy cập nào.

Về hiện thực, pipeline ETL bằng Spark dựng lại toàn bộ kho trong khoảng 5 phút và nạp vào ba
hệ lưu trữ nhất quán với nhau. Airflow điều phối pipeline theo lịch, với cổng chất lượng ngăn
dữ liệu sai đến tay người dùng. Bộ 289 bài kiểm thử bảo đảm số liệu của kho khớp với dữ liệu
nguồn và với đặc tả nghiệp vụ; 20 câu hỏi mẫu cho cùng đáp án trên cả ba hệ lưu trữ. Các kỹ
thuật tối ưu giảm lượng dữ liệu quét trên BigQuery từ 16,8 đến 330 lần.

Về nghiệp vụ, hệ thống trả lời được cả 8 yêu cầu. Kết quả nổi bật là quy tắc cảnh báo sớm:
chỉ với dữ liệu đến hết tuần 4, 80,6\% sinh viên trong danh sách cảnh báo thực sự rớt hoặc rút
khi đợt học kết thúc, gần gấp đôi tỷ lệ nền. Thiết kế bảng ảnh chụp theo tuần không chứa kết
quả cuối kỳ bảo đảm kết quả này không bị rò rỉ thông tin từ tương lai.

\subsection{Hướng phát triển}
\label{sec:huong-phat-trien}

\textbf{Nạp tăng dần và bắt thay đổi dữ liệu (CDC).} Hiện tại kho được dựng lại toàn bộ mỗi
lần, vì nguồn là bản xuất tĩnh và gói Sandbox của BigQuery không cho dùng lệnh
\texttt{MERGE}. Nếu nguồn là một cơ sở dữ liệu vận hành đang thay đổi liên tục, kiến trúc đề
xuất như sau. Debezium đọc nhật ký giao dịch (WAL) của PostgreSQL nguồn và phát mỗi thay đổi
thêm, sửa, xóa thành một thông điệp vào Kafka, mỗi bảng nguồn một topic. Spark Structured
Streaming đọc các topic này, giữ lại thay đổi mới nhất của mỗi khóa trong từng lô nhỏ, rồi ghi
vào một bảng staging chỉ ghi thêm trên BigQuery. Airflow định kỳ chạy \texttt{MERGE} từ bảng
staging vào kho: dimension được cập nhật thuộc tính (SCD loại 1, hoặc loại 2 với thuộc tính
cần lưu lịch sử), fact được thêm dòng mới, còn dòng bị xóa ở nguồn được đánh dấu bằng cờ thay
vì xóa khỏi kho; bảng tiến độ theo tuần chỉ được tính lại cho các lượt đăng ký có thay đổi.

Một bước trung gian nhẹ hơn là nạp tăng dần theo đợt học: mỗi lần chỉ dựng và nạp lại các đợt
có dữ liệu mới. Spark hỗ trợ ghi đè từng phân vùng, còn BigQuery cho phép load job ghi đè đúng
một phân vùng của bảng. Cách này phù hợp với OULAD vì dữ liệu vốn đã chia theo đợt học, nhưng
chưa được thử nghiệm trong project.

\textbf{Đưa các tối ưu vào luồng chính.} Phân vùng, gom cụm và chỉ mục cần được áp dụng ngay
khi nạp kho thay vì chỉ thử nghiệm trên bản sao, và bước làm mới data mart cần được thêm vào
DAG sau bước nạp BigQuery.

\textbf{Tiến tới môi trường production.} Bảng~\ref{tab:production} so sánh hiện trạng của
project với yêu cầu của một hệ thống vận hành thật.

\begin{table}[htbp]
\centering
\small
\begin{tabularx}{\textwidth}{|p{2.6cm}|X|X|}
\hline
\textbf{Thành phần} & \textbf{Yêu cầu khi vận hành thật} & \textbf{Hiện trạng} \\ \hline
Spark & Chạy trên cụm (Dataproc, Databricks, Kubernetes) & Một máy, chế độ local \\ \hline
BigQuery & Bật tính phí để dùng lệnh sửa dữ liệu, ghi stream độ trễ thấp, bảng không tự hết hạn & Gói Sandbox: chỉ load job, bảng tự xóa sau 60 ngày \\ \hline
Airflow & Phiên bản đang được hỗ trợ (3.x), executor chạy trên nhiều máy, cảnh báo khi DAG chạy quá lâu & Phiên bản 2.10.5, \texttt{LocalExecutor}; đã có báo lỗi khi task thất bại \\ \hline
CI/CD & Tự động kiểm thử, dựng image và triển khai mỗi lần cập nhật mã nguồn & Chạy kiểm thử thủ công \\ \hline
Giám sát & Cảnh báo dữ liệu cũ quá hạn, số dòng bất thường; truy vết nguồn gốc dữ liệu (lineage) & Log và báo lỗi của Airflow; cổng chất lượng chặn số dòng sai \\ \hline
Bảo mật & Quản lý khóa bí mật tập trung, phân quyền theo vai trò, che dữ liệu cá nhân & Khóa lưu trong file cục bộ; mọi thành viên cùng một quyền \\ \hline
\end{tabularx}
\caption{Hiện trạng của project so với yêu cầu vận hành thật}
\label{tab:production}
\end{table}
% TODO (bổ sung sau): thêm dòng Kafka (1 broker -> ít nhất 3 broker, replication factor 3) khi viết phần stream.

\subsection{Mã nguồn}
% TODO: link repo GitHub. Ghi rõ trong repo có: hợp đồng bảng và cột (docs/contract.md),
%       20 câu hỏi mẫu (sql/questions/), hướng dẫn chạy lại (README.md),
%       bảng truy vết chuẩn đầu ra đầy đủ (docs/lo_traceability.md).
Toàn bộ mã nguồn, hợp đồng dữ liệu, SQL của 20 câu hỏi mẫu và hướng dẫn dựng lại hệ thống
được lưu tại kho mã nguồn của nhóm: % TODO: \url{...}
